\hspace{3mm}

\centerline{\bf \Huge Страшное слово на букву "K"}

\begin{flushright}
    \scriptsize{}
\end{flushright}

\problem{1}
$\bigl[\textbf{Финал ВсОШ, 1999, 10.8}\bigr]$
В некоторой группе из 12 человек среди каждых девяти найдутся пять попарно знакомых. Докажите, что в этой группе найдутся шесть попарно знакомых.

\problem{2}
$\bigl[\textbf{Финал ВсОШ, 2001, 11.7}\bigr]$
В стране 2001 город, некоторые пары городов соединены дорогами, причём из каждого города выходит хотя бы одна дорога и нет города, соединённого дорогами со всеми остальными. Назовём множество городов D доминирующим, если каждый не входящий в D город соединён дорогой с одним из городов множества D. Известно, что в каждом доминирующем множестве хотя бы k городов. Докажите, что страну можно разбить на  2001 – k  республик так, что никакие два города из одной республики не будут соединены дорогой.

\problem{3}
$\bigl[\textbf{Финал ВсОШ, 2012, 9.2}\bigr]$
На окружности отмечены 2012 точек, делящих её на равные дуги. Из них выбрали k точек и построили выпуклый k-угольник с вершинами
в выбранных точках. При каком наибольшем k могло оказаться, что у этого многоугольника нет параллельных сторон?

\problem{4}
$\bigl[\textbf{Задача Эйлера о шляпах}\bigr]$
Войдя в ресторан, n гостей оставили швейцару свои шляпы, а на выходе получили их обратно. Швейцар раздал шляпы случайным образом. Сколько существует вариантов, при которых каждый гость получит чужую шляпу? Какова вероятность того, что каждый взял не свою шляпу?
